% GENERATED FILE. Edit canonical lessons, then run npm run generate:books.
% canonical-source-hash: fnv1a64:9c184abc3e26d0ac
% canonical-lessons: JA-C244-shiriai, JA-C244-deau, JA-C244-wakareru, JA-C244-mukaeru, JA-C244-kekkon

\chapter{Meeting and Parting}
\label{ch:ja-meeting-and-parting}
\hlchaptermodality{244}

\begin{chapteropening}
\textbf{By the end of this chapter:} \emph{I can ask and answer questions about meeting people and parting.}

Complete the last lesson of chapter 244: I can ask and answer questions about meeting people and parting.
\end{chapteropening}

\section[\texorpdfstring{\ja{しりあい} (shiriai)}{しりあい (shiriai)}]{\ja{しりあい} (shiriai) — an acquaintance}

\label{lesson:JA-C244-shiriai}



\begin{quote}
Before the new one: say the Japanese for siblings, then the Japanese for to bring up.
\end{quote}

\subsection*{You'll want to know: \ja{しりあい}}
\textbf{\ja{しりあい}} — \emph{shiriai} — \textquotedblleft{}an acquaintance\textquotedblright{}.

\begin{quote}
\textbf{\ja{かれは} \ja{ちちの} \ja{しりあいですか。}} — \emph{kare wa chichi no shiriai desu ka} — \textquotedblleft{}Is he an acquaintance of my father's?\textquotedblright{}
\end{quote}

\begin{cousinweb}
\textbf{\ja{しりあい}} (\emph{shiriai}) is a native Japanese compound (wago), inherited rather than borrowed: it is built from \emph{shiri}, from \emph{shiru}, \textquotedblleft{}to know\textquotedblright{}, and \emph{ai}, \textquotedblleft{}mutual\textquotedblright{}, from \emph{au} (JMdict, KANJIDIC2).
\end{cousinweb}

\begin{grammarlens}[title={What you've built}]
One more word for everyday life at A2.
\end{grammarlens}

\subsection*{Your turn}
\begin{itemize}
\raggedright
  \item \emph{Say it:} \emph{shiriai}
  \item \emph{Say it:} \emph{shiriai}, once more
  \item \emph{Say it:} \emph{kare wa chichi no shiriai desu ka} — Is he an acquaintance of my father's?
  \item \emph{Recall:} say the Japanese for siblings, then the Japanese for to bring up, then say \emph{shiriai} again
\end{itemize}

\begin{tcolorbox}[breakable,colback=teal!4,colframe=teal!35!black,title={Before you move on}]
What does \ja{しりあい} mean? (\textquotedblleft{}An acquaintance\textquotedblright{}.) Say it once more.
\end{tcolorbox}
\section[\texorpdfstring{\ja{であう} (deau)}{であう (deau)}]{\ja{であう} (deau) — to come across; to meet by chance}

\label{lesson:JA-C244-deau}



\begin{quote}
Before the new one: say the Japanese for to bring up, then the Japanese for an acquaintance.
\end{quote}

\subsection*{You'll want to know: \ja{であう}}
\textbf{\ja{であう}} — \emph{deau} — \textquotedblleft{}to come across; to meet by chance\textquotedblright{}.

\begin{quote}
\textbf{\ja{まちで} \ja{よく} \ja{しりあいに} \ja{であいますか。}} — \emph{machi de yoku shiriai ni deaimasu ka} — \textquotedblleft{}Do you often run into acquaintances in town?\textquotedblright{}
\end{quote}

Polite \textbf{\ja{であいます}}, \emph{deaimasu}; a godan verb.

\begin{cousinweb}
\textbf{\ja{であう}} (\emph{deau}) is a native Japanese compound (wago), inherited rather than borrowed: it is built from \emph{de}, from \emph{deru}, \textquotedblleft{}to come out\textquotedblright{}, and \emph{au}, \textquotedblleft{}to meet\textquotedblright{} (JMdict, KANJIDIC2).
\end{cousinweb}

\begin{grammarlens}[title={What you've built}]
One more word for everyday life at A2.
\end{grammarlens}

\subsection*{Your turn}
\begin{itemize}
\raggedright
  \item \emph{Say it:} \emph{deau}
  \item \emph{Say it:} \emph{deau}, once more
  \item \emph{Say it:} \emph{machi de yoku shiriai ni deaimasu ka} — Do you often run into acquaintances in town?
  \item \emph{Recall:} say the Japanese for to bring up, then the Japanese for an acquaintance, then say \emph{deau} again
\end{itemize}

\begin{tcolorbox}[breakable,colback=teal!4,colframe=teal!35!black,title={Before you move on}]
What does \ja{であう} mean? (\textquotedblleft{}To come across; to meet by chance\textquotedblright{}.) Say it once more.
\end{tcolorbox}
\section[\texorpdfstring{\ja{わかれる} (wakareru)}{わかれる (wakareru)}]{\ja{わかれる} (wakareru) — to part; to break up}

\label{lesson:JA-C244-wakareru}



\begin{quote}
Before the new one: say the Japanese for an acquaintance, then the Japanese for to come across.
\end{quote}

\subsection*{You'll want to know: \ja{わかれる}}
\textbf{\ja{わかれる}} — \emph{wakareru} — \textquotedblleft{}to part; to break up\textquotedblright{}.

\begin{quote}
\textbf{\ja{あねは} \ja{かれと} \ja{わかれません。}} — \emph{ane wa kare to wakaremasen} — \textquotedblleft{}My elder sister is not breaking up with him.\textquotedblright{}
\end{quote}

Polite \textbf{\ja{わかれます}}, \emph{wakaremasu}; an ichidan verb.

\begin{cousinweb}
\textbf{\ja{わかれる}} (\emph{wakareru}) is a native Japanese word (wago), inherited rather than borrowed from Chinese; its kanji was attached later as a kun'yomi (JMdict, KANJIDIC2).
\end{cousinweb}

\begin{grammarlens}[title={What you've built}]
One more word for everyday life at A2.
\end{grammarlens}

\subsection*{Your turn}
\begin{itemize}
\raggedright
  \item \emph{Say it:} \emph{wakareru}
  \item \emph{Say it:} \emph{wakareru}, once more
  \item \emph{Say it:} \emph{ane wa kare to wakaremasen} — My elder sister is not breaking up with him.
  \item \emph{Recall:} say the Japanese for an acquaintance, then the Japanese for to come across, then say \emph{wakareru} again
\end{itemize}

\begin{tcolorbox}[breakable,colback=teal!4,colframe=teal!35!black,title={Before you move on}]
What does \ja{わかれる} mean? (\textquotedblleft{}To part; to break up\textquotedblright{}.) Say it once more.
\end{tcolorbox}
\section[\texorpdfstring{\ja{むかえる} (mukaeru)}{むかえる (mukaeru)}]{\ja{むかえる} (mukaeru) — to welcome; to go to meet}

\label{lesson:JA-C244-mukaeru}



\begin{quote}
Before the new one: say the Japanese for to come across, then the Japanese for to part.
\end{quote}

\subsection*{You'll want to know: \ja{むかえる}}
\textbf{\ja{むかえる}} — \emph{mukaeru} — \textquotedblleft{}to welcome; to go to meet\textquotedblright{}.

\begin{quote}
\textbf{\ja{ははは} \ja{えきで} \ja{ともだちを} \ja{むかえますか。}} — \emph{haha wa eki de tomodachi o mukaemasu ka} — \textquotedblleft{}Is my mother meeting her friend at the station?\textquotedblright{}
\end{quote}

Polite \textbf{\ja{むかえます}}, \emph{mukaemasu}; an ichidan verb.

\begin{cousinweb}
\textbf{\ja{むかえる}} (\emph{mukaeru}) is a native Japanese word (wago), inherited rather than borrowed from Chinese; its kanji was attached later as a kun'yomi (JMdict, KANJIDIC2).
\end{cousinweb}

\begin{grammarlens}[title={What you've built}]
One more word for everyday life at A2.
\end{grammarlens}

\subsection*{Your turn}
\begin{itemize}
\raggedright
  \item \emph{Say it:} \emph{mukaeru}
  \item \emph{Say it:} \emph{mukaeru}, once more
  \item \emph{Say it:} \emph{haha wa eki de tomodachi o mukaemasu ka} — Is my mother meeting her friend at the station?
  \item \emph{Recall:} say the Japanese for to come across, then the Japanese for to part, then say \emph{mukaeru} again
\end{itemize}

\begin{tcolorbox}[breakable,colback=teal!4,colframe=teal!35!black,title={Before you move on}]
What does \ja{むかえる} mean? (\textquotedblleft{}To welcome; to go to meet\textquotedblright{}.) Say it once more.
\end{tcolorbox}
\section[\texorpdfstring{\ja{けっこん} (kekkon)}{けっこん (kekkon)}]{\ja{けっこん} (kekkon) — marriage; a wedding}

\label{lesson:JA-C244-kekkon}



\begin{quote}
Before the new one: say the Japanese for to part, then the Japanese for to welcome.
\end{quote}

\subsection*{You'll want to know: \ja{けっこん}}
\textbf{\ja{けっこん}} — \emph{kekkon} — \textquotedblleft{}marriage; a wedding\textquotedblright{}.

\begin{quote}
\textbf{\ja{あには} \ja{らいねん} \ja{けっこんしますか。}} — \emph{ani wa rainen kekkon shimasu ka} — \textquotedblleft{}Is my elder brother getting married next year?\textquotedblright{}
\end{quote}

\begin{cousinweb}
\textbf{\ja{けっこん}} (\emph{kekkon}) is Sino-Japanese (kango): two Chinese characters read with their on'yomi, readings borrowed from Chinese. They are \emph{ketsu}, \textquotedblleft{}tie\textquotedblright{} (Middle Chinese \emph{ket}) and \emph{kon}, \textquotedblleft{}marriage\textquotedblright{} (\emph{xwon}) (KANJIDIC2; Middle Chinese in Baxter's transcription).
\end{cousinweb}

\begin{grammarlens}[title={What you've built}]
One more word for everyday life at A2.
\end{grammarlens}

\subsection*{Your turn}
\begin{itemize}
\raggedright
  \item \emph{Say it:} \emph{kekkon}
  \item \emph{Say it:} \emph{kekkon}, once more
  \item \emph{Say it:} \emph{ani wa rainen kekkon shimasu ka} — Is my elder brother getting married next year?
  \item \emph{Recall:} say the Japanese for to part, then the Japanese for to welcome, then say \emph{kekkon} again
\end{itemize}

\begin{tcolorbox}[breakable,colback=teal!4,colframe=teal!35!black,title={Before you move on}]
What does \ja{けっこん} mean? (\textquotedblleft{}Marriage; a wedding\textquotedblright{}.) Say it once more.
\end{tcolorbox}
